\documentclass[11pt,letterpaper]{article}

\usepackage[margin=1in]{geometry}
\usepackage{setspace}
\usepackage{booktabs}
\usepackage{hyperref}
\usepackage{enumitem}
\usepackage{titlesec}
\usepackage{xcolor}

\hypersetup{
  colorlinks=true,
  linkcolor=black,
  urlcolor=blue,
  citecolor=black
}

\setstretch{1.03}
\setlist[itemize]{leftmargin=1.4em,itemsep=0.25em,topsep=0.3em}
\setlist[enumerate]{leftmargin=1.4em,itemsep=0.25em,topsep=0.3em}
\titleformat{\section}{\large\bfseries}{\thesection}{0.6em}{}
\titleformat{\subsection}{\normalsize\bfseries}{\thesubsection}{0.6em}{}

\title{\textbf{CAI6734 Week~10 Individual Progress Report}\\[0.35em]
\large Testing in Imagination:\\
World-Model Planning for Language Agents\\
in Sequential Clinical Diagnosis}
\author{Daniel Escalante\\[0.4em]
\small Course: CAI6734 (University of Florida)\\
Advisor: Xuefeng Liu}
\date{October 30, 2026}

\begin{document}
\maketitle
\thispagestyle{empty}

\noindent
This report covers my individual work for Week~10 (October 28 -- November 3, 2026), written through October~30. The main progress is a tested LA-CDM replication pipeline (upstream \texttt{dharouni/LA-CDM} at \texttt{3f435a10}), ready for MIMIC-IV-Ext CDM v1.1 once my access is approved. I also reviewed Edwin's world-model study and finished my dataset-access steps. Pipeline numbers come from \textbf{invented synthetic data} and are checks, not results. No PHI, no made-up metrics.

\section{Progress this week}

\subsection{Replication pipeline built and merged}
\begin{itemize}
  \item I merged \href{https://github.com/scalinity/testing-in-imagination/pull/2}{PR~\#2} into \texttt{main} on October~30 at 12:44~PM ET (merge commit \texttt{41b9520}). It contains 12 commits, all dated October~30 (\texttt{70f7457} through \texttt{c863c6a}), and changes 46 files (+12,722 / $-$6 lines), all inside \texttt{replication-of-original/daniel-escalante/}.
  \item The pipeline has seven stages: S0 environment; S1 raw-schema validation and splits; S2 config wiring; S3 zero-shot baseline; S4 training of the primary arm plus three ablations; S5 once-only final-test evaluation; and S6 reporting. Each stage writes a resumable manifest. S6 computes paired bootstrap statistics and writes an automatic \texttt{REPORT.md}. SLURM launch files for HiPerGator are included.
  \item The backbone is \texttt{nvidia/Llama-3.1-Nemotron-Nano-8B-v1}, pinned at revision \texttt{54641c16} and used as both the policy and the HPI summarizer (documented deviations D-5 and D-3). The primary arm keeps the shipped \texttt{incorrect\_reward = 0}; the $-1$ reward runs as a separate ablation arm (D-1). Details are in the folder's \texttt{README.md} and \texttt{PLAN.md}.
\end{itemize}

\subsection{Verification on synthetic data}
\begin{itemize}
  \item \textbf{Tests:} 242 tests passed, including all 83 existing tests in \texttt{lacdm/tests} (PR~\#2 description).
  \item \textbf{End-to-end run:} a full synthetic S0--S6 run completed on 2,400 invented admissions with the 957/648/257/538 class mix, missing tests, and repeated subjects. All five arms (primary, zero-shot, decision-only, no-cost, incorrect~$=-1$) ran under both split protocols (the seed-269 parity split and a subject-grouped split), with 10,000 paired bootstrap resamples.
  \item \textbf{Resume check:} all seven stages resumed with verified artifact hashes and no recomputed episodes. The aggregate outputs (\texttt{REPORT.md}, \texttt{metrics.json} under \texttt{results/synthetic/}) carry an explicit ``SYNTHETIC --- invented data, mock updates are not learning'' banner.
\end{itemize}

\subsection{Data-leakage audit and a ledger/code discrepancy}
\begin{itemize}
  \item The seed-269 parity split has 0 admission-ID overlap across train/val/test, but 175 subjects appear in more than one split (80 of them shared between train and test). The subject-grouped split has 0 subject overlap, and the full arm matrix is reported for it as a sensitivity analysis (``Leakage and split sensitivity'' tables in \texttt{REPORT.md}).
  \item The shared ledger's finding D-4 (\texttt{REPRODUCTION.md}) says there is no canonical split seed, but the pinned upstream splitter hardcodes \texttt{RANDOM\_STATE = 269}. The pipeline uses 269 for parity and records the discrepancy in the README and report rather than hiding it.
\end{itemize}

\subsection{Documented gaps, compatibility repairs, and the day-one procedure}
\begin{itemize}
  \item Every manifest and report lists the local compatibility repairs to the upstream training path, for example per-patient validation batches (so one patient's label is not reused), hypothesis SFT batches kept separate from the eight decision rollouts, and one phase per optimizer update. Known upstream issues F-1 to F-4 and the shipped primary reward are preserved.
  \item The ReAct baseline (B2) is recorded as a gap, because the pinned upstream ships no ReAct implementation. I did not invent a substitute baseline.
  \item The README now has a ``Day one: credentialed MIMIC data on HiPerGator'' section: set \texttt{.env} (data and work directories outside the repository, SLURM account and partition), validate the schema before requesting GPU time, do a SLURM dry run, then submit. Only aggregate metrics and \texttt{REPORT.md} are ever exported to the repository. Synthetic run-log rows dated 2026-10-30 are labeled SYNTHETIC, no weight updates.
\end{itemize}

\subsection{Peer review of Edwin's PR~\#1 (synthetic world-model study)}
\begin{itemize}
  \item On October~29 (about 9:18~PM ET) I posted a comment review with 15 line comments on Edwin's \href{https://github.com/scalinity/testing-in-imagination/pull/1}{PR~\#1} (commit dated October~28). Main points, offered as suggestions:
  \begin{enumerate}[label=(\alph*),itemsep=0.05em,topsep=0.1em]
    \item The ``held-out'' log-likelihoods are not held out: the environment's RNG ignores the split, so validation replays the first 500 training patients.
    \item The imagination planner charges full cost for tests that turn out to be missing, so the Oracle arm is not a true upper bound.
    \item With 3 unpaired seeds, the 0.024 world-model-over-model-free gap is suggestive (Welch $p \approx 0.07$); I suggested paired per-seed differences over 10--20 seeds.
    \item The model-free baseline is a 16-state tabular one-step Q computed through the learned model.
    \item History and informative missingness can be added to the synthetic environment now rather than waiting for MIMIC.
  \end{enumerate}
\end{itemize}

\subsection{Dataset access and team coordination}
\begin{itemize}
  \item I completed the two dataset steps listed as next steps in Week~9: I submitted my CITI ``Data or Specimens Only Research'' completion report to the PhysioNet Training page (October~27, after the Week~9 report was filed) and signed the MIMIC-IV-Ext CDM v1.1 data use agreement (October~29), requesting dataset access, as I told Hugo and Edwin by email on October~29. Approval is still pending (PhysioNet lists the training report as under review), and no data has been downloaded.
  \item On October~29 (9:56~PM ET) I confirmed to Hugo that his October~27 screenshots show access to the dataset itself, not just a PhysioNet account, and asked him to start the replication in his own folder. At 9:59~PM I asked Hugo and Edwin how long their dataset approval took. Hugo replied at 10:16~PM that it took about a business week; Edwin has not replied yet.
\end{itemize}

\section{Setbacks and how I addressed them}
\begin{itemize}
  \item \textbf{Dataset approval still pending.} I still cannot run anything on real data. Based on Hugo's experience, approval should take roughly a business week from the request. Instead of waiting, I built and tested the pipeline on synthetic data and wrote a day-one procedure. S1 on real v1.1 files has not run; the HiPerGator \texttt{.env} paths and SLURM settings are not set; no CUDA install (incl.\ the \texttt{xgrammar} source build) or SLURM job has run.
  \item \textbf{Subject leakage in the parity split, and the D-4 seed mismatch.} The seed-269 parity split is admission-disjoint but shares 175 subjects across splits (80 between train and test). I kept it for parity with the paper and added a subject-grouped split (0 subject overlap) that runs the full arm matrix as a sensitivity analysis. Because the ledger (D-4) says there is no canonical seed while the pinned code hardcodes 269, I used 269 and recorded the discrepancy openly.
  \item \textbf{ReAct baseline (B2) cannot be assessed} at the pinned upstream; it is recorded as a gap, not replaced with an invented baseline.
  \item \textbf{Issues found in the PR~\#1 review} (items a--e above) were shared with Edwin as suggestions, with an offer to pair on fixes.
\end{itemize}

\section{Next steps}
\begin{enumerate}
  \item Once access is approved, download MIMIC-IV-Ext CDM v1.1 to private HiPerGator storage (never the repository) and set the \texttt{.env} paths and SLURM settings.
  \item Run schema validation, then the SLURM dry run, then submit S0--S6 on the real data.
  \item Log the first real-data run in the README run log (aggregate metrics only).
  \item Follow up with Edwin (approval time and the PR~\#1 review findings), and support Hugo in starting his replication folder; both folders still hold only starter READMEs.
\end{enumerate}

\section*{Notes}
The synthetic run (mock policy, fake HPI summaries, no weight updates) shows only that the pipeline executes, resumes, and reports; not learning, clinical validity, or paper reproduction. Real v1.1 files, CUDA, training, and scheduler execution are unverified. No Table~1 accuracy, F1, or cost figures are reported.

\end{document}
